\documentclass[12pt]{article}
\usepackage{graphicx}
\usepackage[margin=1in]{geometry}
\usepackage{caption}
\usepackage{xcolor}
\usepackage{float}
\usepackage{amsmath}
\usepackage{array}
\usepackage{multicol}
\usepackage{listings}
\usepackage{hyperref}
\usepackage{url}

\setlength{\columnsep}{1cm}

\definecolor{codegray}{gray}{0.95}
\lstset{
    backgroundcolor=\color{codegray},
    basicstyle=\ttfamily\small,
    columns=fullflexible,
    frame=single,
    breaklines=true,
    showstringspaces=false
}

\begin{document}

% Title Block
\begin{figure}[H]
    \begin{minipage}{0.45\textwidth}
        \includegraphics[width=\textwidth]{WhatsApp Image 2025-07-08 at 3.59.18 PM.jpeg} % Replace with image of circuit setup
    \end{minipage}
\hfill
    \begin{minipage}{0.45\textwidth}
        \textbf{Name: Madhu Latha G} \\
        \textbf{ID: cometfwc029} \\
        \textbf{Date: 9th July 2025}
    \end{minipage}
\end{figure}

\begin{center}
    {\LARGE \textbf{\textcolor{cyan}{GATE Question Paper 2010, EE Question Number 55}}}
\end{center}
\vspace{1em}

% Abstract
\noindent\textbf{Abstract} \\[0.5em]
\includegraphics[width=0.99\linewidth]{55.png} \\[0.3em]
\textit{(GATE 2010, Question No. 55 – Thyristor Commutation using LC Circuit simulated on Arduino)} \\

\vspace{1em}
\noindent\textbf{1. Components Used}
\begin{table}[H]
\small
\centering
\begin{tabular}{|p{4.2cm}|c|}
\hline
\textbf{Component} & \textbf{Qty} \\
\hline
Arduino UNO & 1 \\
Breadboard & 1 \\
LED & 1 \\
220$\Omega$ Resistor & 1 \\
Jumper Wires & 5 \\
USB Cable (Type B) & 1 \\
Computer with Arduino IDE & 1 \\
\hline
\end{tabular}
\caption*{Table 1: List of components used}
\end{table}

\vspace{1em}
\noindent\textbf{2. Setup and Connections}
\begin{enumerate}
    \item Connect digital pin 9 to a 220$\Omega$ resistor.
    \item Connect resistor to LED anode; LED cathode to GND.
    \item Connect Arduino UNO to PC via USB.
    \item Upload sketch using Arduino IDE.
\end{enumerate}

\vspace{1em}
\noindent\textbf{3. Simulation Logic}
\begin{itemize}
    \item The LC circuit causes the current through a thyristor to follow $i(t) = 5\cos(5000t)$.
    \item At $i(t) = 0$, the thyristor turns OFF ($t \approx 314\,\mu s$).
    \item We simulate this behavior with a cosine fading LED using PWM output.
\end{itemize}

\vspace{1em}
\noindent\textbf{5. Observations and Output}
\begin{itemize}
    \item LED starts at full brightness (5A equivalent).
    \item Brightness fades smoothly to OFF in 314\,$\mu$s.
    \item Simulates the current reversal and turn-off of the thyristor.
\end{itemize}

\vspace{1em}
\noindent\textbf{6. Result}
The Arduino successfully simulates the LC commutation waveform. The LED visually represents thyristor current from 5 A to 0 A, confirming the correct turn-off time ($314\,\mu$s).

\vspace{1em}
\noindent\textbf{7. Source Code Repository} \\
\url{https://github.com/madhuu34/Madhulathafwc/blob/main/Hardware/esp32/hardware_code.py}

\end{document}

